\section{Related Work}
\label{sec:related}

\subparagraph*{Verified analyzers.}
Verasco is a Coq-verified analyzer for most of C99, excluding recursion and
dynamic allocation~\cite{jourdan15}. Its interpreter iterates over the
structure of C\#minor, reanalyzes a function at every call site, and raises an
alarm at a possible recursive call. It covers a far larger language than
Voblint, so the comparison concerns design: Voblint turns calls into equations
over unknowns indexed by node and context and computes them with a generic
solver. The value analysis of Blazy et al.\ checks the candidate of an
untrusted iterator with a verified checker~\cite{blazy13}; Voblint instead runs
a solver whose partial correctness is proved. Nipkow and Klein develop
abstract interpretation for IMP in Isabelle/HOL~\cite{nipkow12,nipkow14}, and
Voblint's non-relational domains follow their generic interpreter. Their
language has no procedures. Franceschino et al.\ verify an IMP analyzer in F*
that runs in a browser~\cite{franceschino21}, and Cachera and Pichardie extract
a certified denotational interpreter~\cite{cachera10}; neither has procedures.
Michelland et al.\ make soundness proofs reusable by sharing one interpreter
between the concrete and the abstract semantics~\cite{michelland24}. Cachera et al.\ compute a context-insensitive
dataflow analysis with a verified solver for ordinary
inequations~\cite{cachera05}.

\subparagraph*{Contexts in mechanized semantics.}
Dabrowski and Pichardie instrument a concrete semantics with contexts through
a module type whose call-context function computes a callee's context, and
prove a context-sensitive points-to analysis sound in Coq~\cite{dabrowski09}.
This is the closest precedent for Voblint's context-indexed semantics. Their
analysis is specified by constraints that no solver computes. Voblint's
contexts are admitted by a relation, the totality condition makes the indexed
collection exhaustive, and the indexing is connected to an executable verified
solver. Trace partitioning indexes sets of traces by control
history~\cite{rival07}; Rival and Mauborgne allow coverings in which one trace
has several indices and present calling contexts as one such indexing. The
activation collecting semantics is a covering of this relational kind. Among
Isabelle developments, interprocedural analyses have been
verified~\cite{lammich07afp,wasserrab09afp}, but none of these entries claims an extracted analyzer or represents an activation with its caller chain.

\subparagraph*{Verified solvers.}
Hofmann et al.\ verify a local generic solver in Coq~\cite{hofmann10}, and de
Vilhena et al.\ verify one that discovers dependencies dynamically in
Iris~\cite{devilhena20}. Stade et al.\ prove the top-down solver partially
correct in Isabelle/HOL~\cite{stade24}, and Tilscher et al.\ extend it to side
effects with a generic interface for update rules~\cite{tilscher26}. Voblint
vendors the latter and connects its certificate to programs. Termination is
proved for top-down solvers without side effects over a finite type of
unknowns~\cite{tilscher26jar}; Voblint's unknowns pair nodes, of an infinite type, with contexts, so that
result does not apply.

\subparagraph*{Goblint and local traces.}
Side-effecting constraint systems~\cite{seidl03,apinis12} and the split
between analyses and solvers are Goblint's architecture~\cite{vojdani16,seidl26},
and Voblint's \isaname{dg_spec} follows its \texttt{Spec} interface. The
correspondence is architectural; no theorem transfers to Goblint's OCaml
implementation. Schwarz et al.\ prove Goblint's thread-modular analyses sound
against local traces on paper~\cite{schwarz21,schwarz23}. Activation traces adapt their construction to procedure activations.
Activations do not interfere, so no concurrency result of that work
transfers.
